(1) 旧基态满足 $\langle x\rangle=\langle p\rangle=0$、$\langle p^2\rangle/(2m)=\hbar\omega/4$、$\langle m\omega^2x^2/2\rangle=\hbar\omega/4$。新势相对旧势多出中心偏置的平均能 $m\omega^2x_0^2/2$，故 $\boxed{\langle H_f\rangle=\hbar\omega/2+m\omega^2x_0^2/2}$。

(2) 定义零点长度 $x_{\rm zpf}=\sqrt{\hbar/(2m\omega)}$，新湮灭算符 $a_f=[(x-x_0)/(2x_{\rm zpf})]+ip/(2m\omega x_{\rm zpf})=a_i-x_0/(2x_{\rm zpf})$。旧真空满足 $a_i|0_i\rangle=0$，所以 $a_f|0_i\rangle=\alpha|0_i\rangle$，$\boxed{\alpha=-x_0/(2x_{\rm zpf}),\quad\bar n=|\alpha|^2=m\omega x_0^2/(2\hbar)}$。

(3) 相干态在新Fock基展开为 $|\alpha\rangle=e^{-|\alpha|^2/2}\sum_{n=0}^\infty\alpha^n|n_f\rangle/\sqrt{n!}$。因此 $\boxed{P_n=e^{-\bar n}\bar n^n/n!}$，零激发概率 $P_0=e^{-\bar n}$，粒子数均值与方差都等于 $\bar n$，不同于压缩真空仅出现偶数激发的分布。

(4) 新势阱中Ehrenfest方程 $\dot{\langle x\rangle}=\langle p\rangle/m$，$\dot{\langle p\rangle}=-m\omega^2(\langle x\rangle-x_0)$。初值均为零给 $\boxed{\langle x(t)\rangle=x_0(1-\cos\omega t),\quad\langle p(t)\rangle=m\omega x_0\sin\omega t}$。Gaussian波包宽度始终等于基态宽度，只是中心沿经典轨迹振动。
